% Templat TA TIF UIN Suska Riau -- (c) 2026 Benny Sukma Negara
\chapter{Kajian Pustaka}
\label{bab:kajian}

%==========================================================
\section{Kajian Metode}
\label{sec:kajianmetode}

\subsection{Deskripsi}
Kajian metode berisi kerangka teori yang digunakan di dalam penelitian yang diperoleh dari buku referensi ilmiah, artikel jurnal dan publikasi ilmiah lainnya yang sesuai. Adapun sumber referensi yang disarankan antara lain \parencite{pressman2020software}:
\begin{enumerate}
  \item Buku referensi ilmiah
  \item Artikel/\textit{paper} pada jurnal ilmiah
  \item Buku panduan, undang-undang, dokumen kedinasan dan dokumen resmi.
\end{enumerate}

Adapun susunan isi yang disarankan adalah:
\begin{enumerate}
  \item Deskripsi teori/konsep yang digunakan (beserta referensi).
  \item Deskripsi algoritma/metode yang digunakan (beserta referensi).
\end{enumerate}

\subsection{[Metode/Algoritma yang Digunakan]}
\isi{Uraikan metode lengkap dengan rumus dan penjelasan simbol.}

Contoh penulisan rumus bernomor yang otomatis masuk Daftar Rumus:
\begin{rumus}{Akurasi}
  \text{Akurasi} = \frac{TP + TN}{TP + TN + FP + FN}
  \label{eq:akurasi}
\end{rumus}
dengan $TP$ adalah \textit{true positive}, $TN$ adalah \textit{true negative}, $FP$ adalah \textit{false positive}, dan $FN$ adalah \textit{false negative}. Rujuk rumus dengan Rumus~\eqref{eq:akurasi}.

%==========================================================
\section{Penelitian Terkait}
\label{sec:terkait}

Penelitian terkait berisi deskripsi terkini terkait topik penelitian yang dilakukan, data-data penunjang terkait dengan motivasi penelitian, metodologi penelitian, dan teori/konsep yang digunakan. Ringkasannya disajikan pada Tabel~\ref{tab:terkait}.

\begin{table}[htbp]
\caption{Ringkasan penelitian terkait}
\label{tab:terkait}
\begin{tabularx}{\textwidth}{|c|>{\raggedright\arraybackslash}p{30mm}|>{\raggedright\arraybackslash}X|>{\raggedright\arraybackslash}X|}
\hline
No & Peneliti (Tahun) & Metode dan Hasil & Perbedaan dengan Penelitian Ini\\
\hline
1 & \isi{Nama (tahun)} & \isi{...} & \isi{...}\\
\hline
2 & \isi{Nama (tahun)} & \isi{...} & \isi{...}\\
\hline
\end{tabularx}
\end{table}
